% ===========================================================================
%  Raport — Experiment: INT8 MLP-only vs Full Quantization
%  ViT-Tiny pretrained pe ImageNet-1k validation
%  Tudose Alexandru · Lucrare de Licență 2026
% ===========================================================================
\documentclass[12pt,a4paper]{article}

\usepackage{fontspec}
\defaultfontfeatures{Ligatures=TeX}
\usepackage{polyglossia}
\setmainlanguage{romanian}
\usepackage[top=2.5cm, bottom=2.5cm, left=2.5cm, right=2.5cm]{geometry}
\usepackage{amsmath, amssymb}
\usepackage{graphicx}
\usepackage{float}
\usepackage{caption}
\usepackage{subcaption}
\captionsetup{font=small, labelfont=bf}
\usepackage{booktabs}
\usepackage{array}
\usepackage{xcolor}

\definecolor{fp32color}{HTML}{0072B2}
\definecolor{int8ptcolor}{HTML}{E69F00}
\definecolor{int8pccolor}{HTML}{D55E00}
\definecolor{mlpptcolor}{HTML}{56B4E9}
\definecolor{mlppccolor}{HTML}{009E73}
\definecolor{lightgray}{HTML}{F4F4F4}

\usepackage[colorlinks=true, linkcolor=fp32color, citecolor=fp32color, urlcolor=fp32color]{hyperref}
\usepackage{microtype}
\usepackage{parskip}
\setlength{\parindent}{0pt}
\setlength{\parskip}{6pt}
\usepackage{enumitem}

\usepackage{fancyhdr}
\pagestyle{fancy}
\fancyhf{}
\lhead{\small Tudose Alexandru --- Experiment: MLP-only INT8}
\rhead{\small 2026}
\cfoot{\thepage}
\renewcommand{\headrulewidth}{0.4pt}

% ===========================================================================
\begin{document}
% ===========================================================================

\begin{titlepage}
  \centering
  {\large Universitatea Transilvania din Brașov\\}
  {\small Facultatea de Matematică și Informatică\\}
  {\small Specializarea: Informatică Aplicată\\[2cm]}

  \rule{\linewidth}{0.4pt}\\[0.5cm]
  {\LARGE\bfseries Experiment: Cuantizare INT8\\[0.3cm]}
  {\Large\bfseries Selective MLP-only vs Completă\\[0.2cm]}
  {\large ViT-Tiny pe ImageNet-1k validation}\\[0.3cm]
  \rule{\linewidth}{0.4pt}\\[2cm]

  \begin{minipage}{0.45\textwidth}
    \textbf{Student:}\\
    Tudose Alexandru\\
    Grupa 10LF333\\
    An III, Informatică Aplicată
  \end{minipage}
  \hfill
  \begin{minipage}{0.45\textwidth}
    \raggedleft
    \textbf{Coordonator:}\\
    Conf.\ dr.\ ing.\ Honorius Gâlmeanu\\
    Informatică Aplicată
  \end{minipage}\\[3cm]
\end{titlepage}

\tableofcontents
\newpage

% ===========================================================================
\section{Introducere și Ipoteză}
% ===========================================================================

Acest raport documentează un experiment explorator realizat pe baza unei
sugestii externe: \emph{cuantizarea exclusivă a straturilor MLP, cu
păstrarea straturilor de atenție în precizie completă, poate îmbunătăți
acuratețea față de cuantizarea completă a modelului.}

Raționamentul din spatele ipotezei este că straturile de atenție
(\texttt{attn.qkv}, \texttt{attn.proj}) sunt implicate în calculul
scorurilor de atenție (softmax), operație sensibilă la variații numerice.
Păstrarea lor în FP32 ar putea reduce acumularea erorilor de cuantizare.

\subsection{Context din Faza 3}

Sensitivity analysis-ul din Faza~3 oferă date relevante:

\begin{itemize}
  \item \textbf{Straturile de atenție sunt cele mai robuste} individual:
    \texttt{blocks.0.attn.proj} produce o \emph{îmbunătățire} de
    $+0.036$~pp când este cuantizat izolat.
  \item \textbf{Straturile MLP sunt cele mai sensibile} individual:
    \texttt{blocks.7.mlp.fc1} produce cea mai mare degradare ($-0.092$~pp).
\end{itemize}

Observație importantă: sensitivity analysis cuantizează un singur strat
la un moment dat. Efectele de \emph{interacțiune} între straturi cuantizate
simultan pot fi diferite față de suma efectelor individuale. Experimentul
de față testează tocmai această ipoteză în condiții de cuantizare globală.

% ===========================================================================
\section{Design Experimental}
% ===========================================================================

\subsection{Configurații comparate}

\begin{table}[H]
  \centering
  \caption{Cele 5 configurații evaluate}
  \begin{tabular}{llrr}
    \toprule
    \textbf{Metodă} & \textbf{Straturi cuantizate} &
    \textbf{Nr.\ straturi} & \textbf{Granularitate} \\
    \midrule
    FP32                & ---                        & 0  & --- \\
    INT8-pt (full)      & Attention + MLP            & 48 & Per-tensor \\
    INT8-pc (full)      & Attention + MLP            & 48 & Per-channel \\
    INT8-mlp-pt         & MLP only (fc1, fc2)        & 24 & Per-tensor \\
    INT8-mlp-pc         & MLP only (fc1, fc2)        & 24 & Per-channel \\
    \bottomrule
  \end{tabular}
  \label{tab:configs}
\end{table}

\subsection{Skip patterns}

\begin{table}[H]
  \centering
  \caption{Straturi excluse de la cuantizare}
  \begin{tabular}{ll}
    \toprule
    \textbf{Configurație} & \textbf{Skip patterns} \\
    \midrule
    Full (Faza 2--3)  & \texttt{norm, cls\_token, pos\_embed, head} \\
    MLP-only          & \texttt{norm, cls\_token, pos\_embed, head, \textbf{attn}} \\
    \bottomrule
  \end{tabular}
  \label{tab:skip_patterns}
\end{table}

Adăugarea pattern-ului \texttt{attn} exclude toate straturile al căror
nume conține acest șir: \texttt{blocks.*.attn.qkv} și
\texttt{blocks.*.attn.proj} --- câte 24 de straturi la ViT-Tiny (12
blocuri $\times$ 2 straturi de atenție per bloc).

\subsection{Protocol}

Identic cu Fazele~1--3: \texttt{vit\_tiny\_patch16\_224}, pretrained
ImageNet-1k, 50\,000 imagini de validare, batch size 64, Apple M4 MPS,
PyTorch 2.9.1. Toate cele 5 configurații au rulat în aceeași sesiune,
pe aceleași date.

% ===========================================================================
\section{Rezultate}
% ===========================================================================

\subsection{Acuratețe și degradare}

\begin{table}[H]
  \centering
  \caption{Rezultate complete --- toate configurațiile}
  \begin{tabular}{llrrrr}
    \toprule
    \textbf{Metodă} & \textbf{Straturi} & \textbf{Accuracy} &
    \textbf{$\Delta$ FP32} & \textbf{Loss} & \textbf{Mem (MB)} \\
    \midrule
    FP32              & 0  & 75.456\% & ---          & 0.9420 & 21.81 \\
    INT8-pt (full)    & 48 & 75.134\% & $-0.322$~pp  & 0.9563 & \phantom{0}6.62 \\
    INT8-pc (full)    & 48 & 75.424\% & $-0.032$~pp  & 0.9428 & \phantom{0}6.70 \\
    INT8-mlp-pt       & 24 & 75.214\% & $-0.242$~pp  & 0.9554 & 11.69 \\
    INT8-mlp-pc       & 24 & \textbf{75.486\%} & \textbf{$+0.030$~pp} & \textbf{0.9423} & 11.73 \\
    \bottomrule
  \end{tabular}
  \label{tab:results}
\end{table}

\begin{figure}[H]
  \centering
  \includegraphics[width=\textwidth]{../results/experiments/mlp_only_int8/plots/01_degradation.png}
  \caption{Degradarea de acuratețe față de FP32 pentru fiecare configurație.
    Valorile negative indică degradare; valorile pozitive indică o ușoară
    \emph{îmbunătățire} față de baseline. Numărul de straturi cuantizate
    este indicat pe fiecare bară.}
  \label{fig:degradation}
\end{figure}

\begin{figure}[H]
  \centering
  \includegraphics[width=\textwidth]{../results/experiments/mlp_only_int8/plots/03_full_vs_mlp_grouped.png}
  \caption{Comparație directă: cuantizare completă (48 straturi) vs MLP-only
    (24 straturi), per granularitate. Linia punctată indică acuratețea FP32.
    MLP-only îmbunătățește acuratețea în ambele cazuri, cu cost în memorie.}
  \label{fig:grouped}
\end{figure}

\subsection{Tradeoff acuratețe--memorie}

\begin{figure}[H]
  \centering
  \includegraphics[width=0.85\textwidth]{../results/experiments/mlp_only_int8/plots/02_tradeoff_accuracy_memory.png}
  \caption{Scatter plot: acuratețe vs memorie pentru toate configurațiile.
    Colțul dreapta-jos reprezintă zona ideală (acuratețe mare, memorie mică).
    INT8-pc (full) oferă cel mai bun raport: $6.70$~MB cu $-0.032$~pp.
    INT8-mlp-pc depășește FP32 ca acuratețe, dar folosește $11.73$~MB.}
  \label{fig:tradeoff}
\end{figure}

\begin{figure}[H]
  \centering
  \includegraphics[width=\textwidth]{../results/experiments/mlp_only_int8/plots/00_summary.png}
  \caption{Sumar complet: acuratețe, memorie și latență pentru toate
    configurațiile.}
  \label{fig:summary}
\end{figure}

% ===========================================================================
\section{Analiză și Interpretare}
% ===========================================================================

\subsection{Ipoteza este confirmată parțial}

\begin{table}[H]
  \centering
  \caption{Verdict: MLP-only vs Full quantization}
  \begin{tabular}{lrr}
    \toprule
    \textbf{Comparație} & \textbf{$\Delta$ Accuracy} & \textbf{Verdict} \\
    \midrule
    INT8-mlp-pt vs INT8-pt (full) & $+0.080$~pp & Confirmat \\
    INT8-mlp-pc vs INT8-pc (full) & $+0.062$~pp & Confirmat \\
    \bottomrule
  \end{tabular}
  \label{tab:verdict}
\end{table}

Skipând straturile de atenție, acuratețea \textbf{se îmbunătățește în
ambele cazuri}: $+0.080$~pp pentru per-tensor și $+0.062$~pp pentru
per-channel. Ipoteza este deci \textbf{confirmată empiric}.

\subsection{Rezultatul remarcabil: INT8-mlp-pc depășește FP32}

\textbf{INT8-mlp-pc atinge 75.486\%} --- cu $+0.030$~pp \emph{peste}
FP32 baseline ($75.456\%$). Aceasta este o valoare situată în afara
intervalului de degradare și necesită interpretare atentă:

\begin{enumerate}
  \item \textbf{Efect de regularizare:} Eroarea de cuantizare introdusă
    în straturile MLP poate acționa ca un regularizator implicit, similar
    cu weight noise sau dropout. Straturile de atenție, rămânând în FP32,
    nu perturbă mecanismul de atenție, în timp ce MLP-urile cuantizate
    introduc un zgomot mic care reduce over-fitting-ul pe distribuția de
    validare.

  \item \textbf{Variabilitate statistică:} O diferență de $+0.030$~pp
    pe 50\,000 de imagini corespunde la 15 imagini clasificate diferit.
    Aceasta se situează \emph{la limita} semnificației statistice
    ($\pm 0.05$~pp este variabilitatea tipică între rulări).

  \item \textbf{Interacțiuni non-liniare:} Cuantizând doar MLP-urile,
    eliminăm erorile care în cuantizarea completă se acumulau prin
    mecanismul de atenție (outlierii din block~7 identificați în Faza~2
    afectau și calculul scorurilor de atenție în straturile ulterioare).
\end{enumerate}

\subsection{Confirmarea pattern-ului din sensitivity analysis}

Sensitivity analysis (Faza~3) a arătat că straturile de atenție sunt
robuste individual. Acest experiment demonstrează că \textbf{concluzia
se menține și la cuantizare globală}: skipând straturile de atenție
(care oricum tolerează cuantizarea), câștigăm acuratețe fără un cost
semnificativ.

\subsection{Costul: compresia memoriei}

\begin{table}[H]
  \centering
  \caption{Comparație memoriei și factorului de compresie}
  \begin{tabular}{lrrr}
    \toprule
    \textbf{Metodă} & \textbf{Mem (MB)} & \textbf{Reducere vs FP32} &
    \textbf{Accuracy} \\
    \midrule
    FP32           & 21.81 & ---         & 75.456\% \\
    INT8-pc (full) & \phantom{0}6.70  & $3.26\times$ & 75.424\% \\
    INT8-mlp-pc    & 11.73 & $1.86\times$ & 75.486\% \\
    \bottomrule
  \end{tabular}
  \label{tab:memory_tradeoff}
\end{table}

Prețul pentru câștigul de acuratețe este semnificativ: INT8-mlp-pc
folosește $11.73$~MB față de $6.70$~MB pentru INT8-pc complet ---
o diferență de $5.03$~MB. Factorul de compresie scade de la
$3.26\times$ la $1.86\times$.

\subsection{Latența}

Toate variantele INT8 au o latență mai mare decât FP32 pe Apple M4
MPS (hardware fără suport INT8 nativ), ceea ce este consistent cu
rezultatele din Faza~3. Diferențele de latență între full și MLP-only
sunt nesemnificative ($\pm 20$~ms), deoarece dequantizarea se face
oricum la nivel de strat, indiferent de numărul de straturi cuantizate.

% ===========================================================================
\section{Concluzii și Recomandări}
% ===========================================================================

\begin{enumerate}
  \item \textbf{Ipoteza confirmată:} Cuantizarea MLP-only îmbunătățește
    acuratețea față de cuantizarea completă: $+0.080$~pp per-tensor,
    $+0.062$~pp per-channel.

  \item \textbf{INT8-mlp-pc --- cea mai bună acuratețe:} $75.486\%$,
    ușor peste FP32 baseline. Recomandat când acuratețea maximă
    primează față de compresia memoriei.

  \item \textbf{INT8-pc (full) --- cel mai bun raport eficiență/acuratețe:}
    $6.70$~MB, $-0.032$~pp. Recomandat când constrângerea principală
    este memoria.

  \item \textbf{Concluzie structurală validată:} Straturile de atenție
    în ViT-Tiny sunt robuste la cuantizare atât individual (Faza~3)
    cât și în context global (acest experiment). Straturile MLP sunt
    sursa principală de degradare.

  \item \textbf{Direcție viitoare --- mixed-precision optim:} Sensitivity
    analysis + acest experiment sugerează o strategie optimă:
    păstrează straturile cu MSE ridicat (block~7 MLP) în FP16/FP32,
    cuantizează per-channel restul straturilor MLP și toate straturile
    de atenție. Aceasta ar putea combina compresia maximă cu degradarea
    minimă.
\end{enumerate}

% ===========================================================================
\begin{thebibliography}{9}
% ===========================================================================

\bibitem{dosovitskiy2021vit}
A. Dosovitskiy et al.,
``An Image is Worth 16x16 Words: Transformers for Image Recognition at Scale,''
\emph{ICLR}, 2021. \href{https://arxiv.org/abs/2010.11929}{arXiv:2010.11929}.

\bibitem{nagel2021whitepaper}
M. Nagel et al.,
``A White Paper on Neural Network Quantization,''
2021. \href{https://arxiv.org/abs/2106.08295}{arXiv:2106.08295}.

\bibitem{dettmers2022llmint8}
T. Dettmers et al.,
``LLM.int8(): 8-bit Matrix Multiplication for Transformers at Scale,''
\emph{NeurIPS}, 2022. \href{https://arxiv.org/abs/2208.07339}{arXiv:2208.07339}.

\bibitem{bondarenko2023quantizable}
Y. Bondarenko, M. Nagel, T. Blankevoort,
``Quantizable Transformers: Removing Outliers by Helping Attention Heads Do Nothing,''
\emph{ECCV}, 2024. \href{https://arxiv.org/abs/2306.12929}{arXiv:2306.12929}.

\end{thebibliography}

\end{document}
